\chapter{Creating a MERN Project}

A MERN project is two independent Node.js projects side by side: a
Vite-powered React frontend and an Express backend, talking only over HTTP.

\begin{diagrambox}[Project Layout]
\begin{verbatim}
my-app/
|-- frontend/   (Vite + React, :5173)
|-- backend/    (Express + MongoDB, :5000)
\end{verbatim}
\end{diagrambox}

\section{Scaffolding the Frontend}

\begin{lstlisting}[language=bash]
mkdir my-app && cd my-app
npm create vite@latest frontend -- --template react
cd frontend
npm install
npm run dev
\end{lstlisting}

This gives a working React app with hot reload. We restructure \texttt{src/}
in Chapter 3 and add the Axios API layer in Chapter 7.

\section{Scaffolding the Backend}

\begin{lstlisting}[language=bash]
cd my-app
mkdir backend && cd backend
npm init -y
npm install express mongoose dotenv cors cookie-parser bcryptjs jsonwebtoken
npm install -D nodemon
\end{lstlisting}

\begin{table}[h]
\centering
\begin{tabularx}{\textwidth}{lX}
\toprule
\textbf{Package} & \textbf{Purpose} \\
\midrule
\texttt{express} & HTTP server framework -- routing, middleware \\
\texttt{mongoose} & ODM mapping JS schemas to MongoDB collections \\
\texttt{dotenv} & Loads env vars from \texttt{.env} \\
\texttt{cors} & Enables cross-origin requests from frontend \\
\texttt{cookie-parser} & Parses \texttt{Cookie} header for JWT \\
\texttt{bcryptjs} & Hashes/compares passwords \\
\texttt{jsonwebtoken} & Issues/verifies JWTs \\
\texttt{nodemon} (dev) & Auto-restarts server on file changes \\
\bottomrule
\end{tabularx}
\end{table}

\begin{lstlisting}[language=JavaScript]
// backend/server.js
const express = require("express");
const cors = require("cors");
const cookieParser = require("cookie-parser");
require("dotenv").config();

const app = express();

app.use(cors({ origin: process.env.CLIENT_URL, credentials: true }));
app.use(express.json());
app.use(cookieParser());

app.get("/api/health", (req, res) => {
  res.json({ success: true, message: "API is running" });
});

const PORT = process.env.PORT || 5000;
app.listen(PORT, () => console.log(`Server running on port ${PORT}`));
\end{lstlisting}

\begin{notebox}[NOTE]
\texttt{MONGO\_URI}, \texttt{JWT\_SECRET}, \texttt{PORT}, \texttt{CLIENT\_URL} belong in a
git-ignored \texttt{.env} file. Mongoose connection setup is in Chapter 9.
\end{notebox}

\begin{tipbox}[TIP]
Add \texttt{"dev": "nodemon server.js"} and \texttt{"start": "node server.js"} to
\texttt{package.json} scripts now.
\end{tipbox}
